% Notes of the Montreal seminar, 1970 (cote 129, typed part): archive cote 129, site page 73 (PDF page 74).
% typescript. Draft compared: none.
% AI-transcribed by Claude Opus 5.5 (claude-opus-5-5, adaptive thinking, effort high), Message Batches API, batch mtl129_full (msgbatch_01AGnDKHwy6ATCM7ehUa6263), 2026-11-13; not yet checked by eye.
\scanpage{129/73}
\mspage{\hand{III - 18}}

7. \uline{Suite de composition d'un groupe de Barsotti-Tate}

7.1. Supposons d'abord $S$ réduit à un point $s$. Si $G$ est un
groupe de Barsotti-Tate sur $S$, on pose par définition
\[ G^{o} = \underset{\vec{n}}{\lim}\ G(n)^{o} \qquad\qquad (G(n)^{o} = \text{composante neutre de } G(n)) \]
\[ G^{et} = \underset{\vec{n}}{\lim}\ G(n)/G(n)^{o} \ , \]
et l'on obtient une suite exacte
\[ 0 \to G^{o} \to G \to G^{et} \to 0 \]
On vérifie immédiatement que $G^{o}$ est un groupe de BT, d'où l'on déduit
(5.2) que $G^{et}$ l'est aussi, de sorte que l'on obtient $G$ comme extension
d'un groupe de BT ind-étale par un groupe de BT à fibre connexe. Cet-%
te décomposition est évidemment canonique et fonctorielle en $G$.

7.2. Si la caractéristique résiduelle de $S$ est $p$ et si l'on pose
\[ G^{tor} = \underset{\vec{n}}{\lim}\ G(n)^{mult} \]
où $G(n)^{mult}$ désigne le plus grand sous-groupe de type multiplicatif de
$G(n)$, $G(n)^{mult} \subset G(n)^{o}$, on obtient une filtration
\[ \{0\} \subset G^{tor} \subset G^{o} \subset G \]

Il est immédiat que $G^{tor}$ est un groupe de BT; c'est le plus
grand sous-groupe de BT toroidal de $G$.

On en déduit que $G^{o}/G^{tor}$ est un groupe de BT; c'est un

\margin{73}
